\section{Formal evidence for the three-row theorem}\label{sec:formalization}

The M1--M6 development follows Section \ref{sec:half}, without using
the block classification. A cell has zero-based coordinates in
\texttt{Fin m} $\times$ \texttt{Fin w}. A domino is an unordered
two-element set of adjacent cells, and a tiling is an exact cover.
The definition \texttt{Domino.a n} is the cardinality of the quotient
of actual $3\times2n$ tilings by the four coordinate symmetries.
The coefficient of $X^n$ in \texttt{orbitSeries} is $(a(n):\mathbb Q)$;
neither the target recurrence nor the rational expression defines it.

\begin{center}
\small
\begin{tabular}{cp{0.78\linewidth}}
\toprule
Stage & Mathematical interface to examine\\
\midrule
M1 & Unordered dominoes, exact coverage, terminal holes, the width-zero
guard, coordinate action, fixed sets, and the orbit quotient.\\
M2 & An equivalence between actual partial tilings and state paths,
with two inverse identities, followed by the counting recurrence.\\
M3 & Reflection and half-turn completion/restriction, including actual
internal and crossing dominoes; uniqueness of the $h$-fixed tiling.\\
M4 & Simultaneous induction of support and counts for all widths;
identification of the auxiliary pair with actual cardinalities.\\
M5 & Burnside applied to that action, divisibility, and the empty and
small-board cases; no faithfulness assumption.\\
M6 & Series defined from those counts, denominator identity and
invertibility, the recurrence, and all seven initial values.\\
\bottomrule
\end{tabular}
\end{center}

Write $B=(1-X)(1-4X+X^2)(1-4X^2+X^4)$ and
$N=1-3X-4X^2+10X^3-4X^5+X^6$.
The development first proves $B A=N$, then uses $B(0)=1$ to obtain
$A=NB^{-1}$ in $\mathbb Q[[X]]$. Coefficient comparison at $n+7$
gives the recurrence for every $n\in\mathbb N$, interpreted in
$\mathbb Z$; a separate natural-number version moves all negative
terms to the other side. Coefficients zero through six give the
initial values, and the first recurrence step gives $a(7)=2037$.
The main endpoints, in namespace \texttt{Domino}, are
\texttt{orbit\_series\_paper}, \texttt{orbit\_recurrence\_seven},
\texttt{orbit\_initial\_values}, and \texttt{three\_row\_main\_theorem}.

The recorded October 25, 2026 acceptance uses Lean 4.35.0 and
the following mathlib commit:
\begin{center}
\small\texttt{56a005855ea17e0b0958f73ec801ae140b0a1679}.
\end{center}
It covers 32 mathematical modules and audits 413 explicit declarations,
including auxiliary definitions and instances. The allowed axioms
are \texttt{propext}, \texttt{Classical.choice}, and \texttt{Quot.sound}.
Project sources were rebuilt in a directory without prior project
build products, reusing the checked third-party dependency cache;
this was not a cache-free rebuild of all of mathlib.

This edition deliberately fixes its review scope at M1--M6.
The general-height statements, stabilizers, asymptotics, five-row
certificate, and block-based second proof are not certified by that
six-stage evidence. Later local stages are separate work.
The four coordinate transformations are explicit; a classification
of all Euclidean isometries is not formalized. At zero width the
abstract group acts trivially on the unique empty tiling.
Kernel checking does not decide whether these definitions express
the intended question, or establish bibliographic priority.

The accompanying \path{review/m1-m6/} directory provides a source-pinned
semantic map, a dependency guide, executable declaration probes,
and an unfilled external review form. It distinguishes the compiler
checks from the human judgments still requested. In particular,
the pair in \eqref{eq:two-dimensional} is defined recursively in Lean
and only then proved equal to the geometric counts by
\texttt{profile\_even\_odd}; the existence of its recursion equations
alone would not establish the tiling count.
